% Documentation for the sfs-2487-2024 LaTeX document class.
%
% Copyright (c) 2026 Asko Soukka <asko.soukka@iki.fi>
%
% SPDX-License-Identifier: MIT
%
% This work is distributed under the MIT License; the full license text is
% in the accompanying LICENSE file and at https://opensource.org/license/mit

% Kept in sync with the class by `l3build tag <version>`.
\def\clsversion{2.0}
\def\clsdate{2026/06/11}

\documentclass[11pt]{article}

\usepackage[T1]{fontenc}
\usepackage[utf8]{inputenc}
\usepackage[english]{babel}
\usepackage{mathpazo}
% expansion=false: the T1 typewriter font in this minimal TeX Live is
% bitmap-only, and pdfTeX cannot expand non-scalable fonts.
\usepackage[expansion=false]{microtype}
\usepackage{enumitem}
\usepackage[a4paper,margin=25mm]{geometry}
\usepackage[bookmarksopen,bookmarksnumbered,hidelinks]{hyperref}

\newcommand{\cls}{\textsf{sfs-2487-2024}}
\newcommand{\cmd}[1]{\texttt{\textbackslash #1}}
\newcommand{\opt}[1]{\texttt{[#1]}}

\setlist[description]{font=\normalfont\ttfamily, style=nextline,
  itemsep=.5\baselineskip}

\title{The \cls\ class\\
  \large Finnish standard office documents per SFS 2487:2024}
\author{Asko Soukka\\\texttt{asko.soukka@iki.fi}}
\date{Version \clsversion\ (\clsdate)}

\hypersetup{pdftitle={The sfs-2487-2024 class},
  pdfauthor={Asko Soukka}, pdflang={en}}

\begin{document}
\maketitle

\begin{abstract}
\noindent
The \cls\ class typesets Finnish office documents --- letters, memos,
minutes, quotations and the like (kirjeet, muistiot, pöytäkirjat,
tarjoukset) --- following the Finnish standard \emph{SFS 2487:2024
Asiakirjan asettelu ja metatiedot} (document layout and metadata). The
class lays out the standard's information areas (tietoalueet)
automatically: the basic metadata area 9,2\,cm from the left margin with
\mbox{``1 (2)''} page numbering, body text on the 2,3\,cm column grid,
headings hanging at the left margin, and the metadata block repeated as a
page header from page~2 onward.
\end{abstract}

\tableofcontents

\section{Introduction}

SFS 2487 is the Finnish standard specifying the layout, information areas
and metadata of formal office documents; the 2024 edition is its current
revision. Documents following it share a fixed geometry: a 2,3\,cm column
grid starting 2\,cm from the paper's left edge, a basic metadata area
(document type, date, identifier, confidentiality, page number) starting
9,2\,cm from the left margin, and body text indented one grid column with
headings hanging at the margin.

The \cls\ class implements these rules on top of the standard
\textsf{article} class so that the author writes only content and
metadata; positioning, fonts, spacing, page numbering and the repeated
continuation-page header all follow from the class. The implementation
has been verified against the standard's own model documents (its
appendices Liite~A and~B).

A minimal document:

\begin{verbatim}
\documentclass{sfs-2487-2024}

\doctype{Pöytäkirja}
\date{15.5.2024}
\author{Virve Virtanen}
\title{Asiakaspalautteet ja etusivun uudistaminen}

\begin{document}
\maketitle

\section{Kokouksen avaus}

Puheenjohtaja avasi kokouksen ja toivotti kaikki
tervetulleiksi.

\end{document}
\end{verbatim}

Build with \texttt{latexmk -pdf}, which automatically runs the extra pass
needed for the total page count in the ``1~(2)'' page number.

\section{Class options}

\begin{description}
\item[(default), sans-serif] Helvetica-like sans serif at 11\,pt, bold
  headings numbered \mbox{``1 Otsikko''}, matching the look of the
  standard's own example renderings.
\item[serif] Palatino body font instead of Helvetica.
\item[monospace] Courier typewriter font throughout.
\item[agenda] Meeting-agenda heading numbering with a trailing period,
  \mbox{``1. Otsikko''} --- an established alternative the standard
  permits (clause 6.4.1). Headings stay bold either way, as clause 6.4
  requires bold or a larger font size.
\item[12pt, \ldots] Any other option is passed through to
  \textsf{article}, e.g.\ a larger base font size.
\end{description}

Options combine: \verb|\documentclass[agenda,serif]{sfs-2487-2024}|.

\section{Metadata commands}

Set these in the preamble, before \verb|\begin{document}|. Only
\cmd{doctype}, \cmd{date}, \cmd{author} and \cmd{title} are usually
needed; everything else is optional and omitted from the output when
unset.

\begin{description}
\item[\textbackslash doctype\{Pöytäkirja\}] The document type
  (asiakirjatyyppi), shown bold at the start of the basic metadata area.
\item[\textbackslash date\{15.5.2024\}] The date (päivämäärä). Write as
  \texttt{d.m.yyyy} without leading zeros, per SFS 4175. Also stored as
  the PDF CreationDate and ModDate, so the document properties agree
  with the visible metadata instead of recording the compilation time.
\item[\textbackslash modified\{5.6.2024\}] Overrides the PDF ModDate
  alone, as \texttt{d.m.yyyy}; without it both PDF dates follow
  \cmd{date}. Not shown on the document. Optional.
\item[\textbackslash author\{Virve Virtanen\}] The author (laatija), part
  of the standard's required minimum metadata (clause 5.2); stored in the
  PDF Author property. To show it on the document itself, add a line to
  \cmd{extrametadata}.
\item[\textbackslash docid\{Dnro 123/2024\}] The document identifier
  (asiakirjan yksilöivä tunnus). Optional.
\item[\textbackslash confidentiality\{Luottamuksellinen\}] The
  confidentiality marking (luottamuksellisuus). Optional.
\item[\textbackslash extrametadata\{Hankenumero 123456\textbackslash\textbackslash{}Asiakasnumero 987654\}]
  Extra metadata lines (lisämetatiedot), separated by \verb|\\|, placed
  below the basic metadata. Optional.
\item[\textbackslash logo\{\ldots\}] The organization logo, as any box:
  \verb|\includegraphics{...}| or text such as
  \verb|\textsf{\textbf{Yritys Oy}}|. Placed top-left; the page header
  grows automatically to fit a tall logo. A logo taller than
  \cmd{logomaxheight} (20\,mm by default) is scaled down
  proportionally first --- adjust with
  \verb|\setlength{\logomaxheight}{30mm}|.
\item[\textbackslash recipient\{Oy Yritys Ab\textbackslash\textbackslash{}\ldots\}]
  The recipient information area (vastaanottajan tietoalue), lines
  separated by \verb|\\|, placed at the left margin below the metadata.
  Optional.
\item[\textbackslash subject\{Digiprojekti\}] The subject (aihe), a bold
  line directly above the main title. Optional.
\item[\textbackslash keywords\{palaute, kysely\}] Keywords (asiasanat)
  for the PDF Keywords property; not shown on the document. Optional.
\item[\textbackslash title\{Verkkosivujen uudistaminen\}] The main title
  (pääotsikko), bold and 3\,pt larger than the body text.
\item[\textbackslash contactinfo\{Organisaatio Oy\}\{Katuosoite\textbackslash\textbackslash{}\ldots\}]
  Organization contact details for \cmd{makecontactinfo}; the first
  argument is the bolded organization name, the second the remaining
  lines separated by \verb|\\|.
\end{description}

The metadata is also written into the PDF document properties --- title,
subject, author, keywords, creation/modification dates and document
language \texttt{fi} --- for accessibility and document management.

\section{Body commands}

\subsection{\textbackslash maketitle}

Renders the standard's information areas 1--5 on the first page: logo,
basic metadata (document type, date, identifier, confidentiality, page
number), extra metadata, recipient, and subject + main title. The same
metadata block repeats automatically as a header on every following
page.

\subsection{Headings}

\begin{verbatim}
\section{Kokouksen avaus}      % 1 Kokouksen avaus
\subsection{Tarkennus}         % 1.1 Tarkennus
\subsubsection{Yksityiskohta}  % 1.1.1 Yksityiskohta
\section*{Esityslista}         % unnumbered, still bold
\end{verbatim}

Headings hang at the left margin; body text after them returns to the
2,3\,cm indent. Up to three numbered levels are available, as the
standard recommends; with the \opt{agenda} class option the numbers get
a trailing period (\mbox{``1. Otsikko''}).

\subsection{\textbackslash marginlabel}

An unnumbered, regular-weight label at the left margin with the
following content at the body indent --- for rows like \emph{Aika ja
paikka} or \emph{Osallistujat}:

\begin{verbatim}
\marginlabel{Osallistujat}

Marja Mäkinen, puheenjohtaja\\
Virve Virtanen, sihteeri
\end{verbatim}

\subsection{Signatures}

An electronic signature block (sähköinen allekirjoitus) produces the
standard's lead-in sentence ``Tämä asiakirja on sähköisesti
allekirjoitettu.''\ followed by name/email pairs:

\begin{verbatim}
\begin{esignatures}
  \esignee{Marja Mäkinen, puheenjohtaja}{marja.makinen@yritys.fi}
  \esignee{Virve Virtanen, sihteeri}{virve.virtanen@yritys.fi}
\end{esignatures}
\end{verbatim}

A handwritten signature (omakätinen allekirjoitus) reserves empty space
for signing above the printed ``Name, role'' line, stackable for several
signers:

\begin{verbatim}
\handsignature{Marja Mäkinen, puheenjohtaja}
\handsignature{Virve Virtanen, sihteeri}
\end{verbatim}

\subsection{Attachments, distribution, for information}

\begin{verbatim}
\attachments{Yhteenveto palautteesta\\Ehdotus muutoksista}
\distribution{Digiprojektin ohjausryhmä}
\forinformation{Johtoryhmä}
\end{verbatim}

Each renders its label (\emph{Liitteet}, \emph{Jakelu}, \emph{Tiedoksi})
at the left margin with the items, one per line, at the body indent. Use
them in this order after the signatures, per the standard.

The class's fixed strings are renewable per the usual
\cmd{refname}/\cmd{contentsname} convention, so a document in another
language keeps the SFS 2487 layout: \cmd{attachmentsname},
\cmd{distributionname}, \cmd{forinformationname} for the labels above,
and \cmd{esignaturestext} for the electronic-signature lead-in
sentence, e.g.\ \verb|\renewcommand{\attachmentsname}{Bilagor}|.

\subsection{\textbackslash makecontactinfo}

Prints the organization contact block set with \cmd{contactinfo} at the
left margin --- place it at the end of the document. The standard
requires contact details in the body area, not only in a footer.

\subsection{Tables and figures}

Captions are styled by the class per clauses 6.5.1--6.5.2: left-aligned
at the text indent, never centered, with the label (\emph{Taulukko~1},
\emph{Kuva~1}) separated by a quad. The standard prefers material placed
in the text flow --- use \cmd{captionof} instead of floating
environments, with the table caption \emph{above} the table and the
figure caption next to the figure:

\begin{verbatim}
\captionof{table}{Vastaukset palautekanavittain}
\begin{tabular}{@{}lrr@{}}
  \textbf{Kanava} & \textbf{Vastauksia} & \textbf{Osuus} \\
  ...
\end{tabular}

\includegraphics{kuva}
\captionof{figure}{Logon perusversio}
\end{verbatim}

The floating \texttt{table}/\texttt{figure} environments work too and
get the same caption styling. Mark the table's top row as a header row
with bold text (clause 6.5.1), and remember an alternative text or a
body-text description for images (clause 6.5.2).

\subsection{Lists and page breaks}

Bullet lists are preconfigured (en-dash markers per Finnish convention,
first level flush at the body text edge, tight item spacing, one
paragraph gap around the list). For looser spacing pass
\textsf{enumitem} options:

\begin{verbatim}
\begin{itemize}[itemsep=\parskip]
  \item 20.5.2024 toimituksen kokoustilassa
  \item 27.5.2024 verkossa
\end{itemize}
\end{verbatim}

Use \cmd{clearpage} to force a page break, e.g.\ before the attachments
area.

\section{Accessibility: tagged PDF}

Documents can opt in to tagged PDF --- the structure tree that the
standard's accessibility appendix (liite~D) builds on --- by adding one
line \emph{before} \cmd{documentclass}:

\begin{verbatim}
\DocumentMetadata{lang=fi-FI, pdfversion=2.0, tagging=on}
\documentclass{sfs-2487-2024}
\end{verbatim}

With tagging on, the output gets a full structure tree: the main title
as a \texttt{Title} element, sections as \texttt{H1}--\texttt{H3},
paragraphs, lists, tables, footnotes, the table of contents, link
annotations, and the logo as a \texttt{Figure}. The repeated page-2+
header is marked as an artifact so assistive technology reads the
metadata once. Give the logo an alternative text:

\begin{verbatim}
\logo{\includegraphics[alt={Organisaatio Oy:n logo}]%
  {logo-organisaatio}}
\end{verbatim}

Untagged builds are byte-for-byte unaffected.

\section{Demo files}

The package ships example documents exercising the class API:

\begin{description}
\item[esimerkki-kokouskutsu.tex] A meeting invitation with an agenda,
  using the \opt{agenda} class option.
\item[esimerkki-raportti.tex] A multi-page report with a table of
  contents, a table and a footnote.
\item[esimerkki-kayttoohje.tex] A manual with captioned figures, using
  the \opt{sans-serif} class option.
\item[esimerkki-monospace.tex] A memo using the \opt{monospace} class
  option for a Courier typewriter look.
\item[logo-*.tex] Invented TikZ logos used by the examples.
\end{description}

The project repository additionally provides a Markdown front end (a
pandoc template and Lua filter), Markdown twins of all examples, and the
full user documentation as a website.

\section{License}

Copyright \textcopyright{} 2026 Asko Soukka. The class and its
documentation are distributed under the MIT License; see the LICENSE
file included in the package.

\end{document}
